\documentclass{article}
\usepackage[T1]{fontenc}
\usepackage{lmodern}
\usepackage{graphicx}
\usepackage{amsmath,amssymb}
\usepackage[a4paper,margin=2cm]{geometry}
\usepackage[absolute,overlay]{textpos}
\setlength{\TPHorizModule}{1mm}
\setlength{\TPVertModule}{1mm}
\usepackage[indonesian]{babel}
\usepackage{hyperref}
\setlength{\emergencystretch}{2em}
% unit: Halaman 10

\begin{document}

% === Halaman 10 (llm) ===

\begin{itemize}
    \item Contoh medan:
    \begin{itemize}
        \item medan bilangan bulat, $\langle\mathbb{Z}, +, .\rangle$
        \item medan bilangan riil, $\langle\mathbb{R}, +, .\rangle$
        \item medan bilangan rasional $(p/q)$, $\langle\mathbb{Q}, +, .\rangle$
    \end{itemize}

    \item Sebuah medan disebut berhingga (\textit{finite field}) jika himpunannya memiliki jumlah elemen yang berhingga.
    
    Jika jumlah elemen di dalam himpunan adalah $n$, maka notasinya $F_n$
    
    Contoh: $F_2$ adalah medan dengan elemen 0 dan 1

    \item Medan berhingga sering dinamakan juga \textcolor{red}{Galois Field}, untuk menghormati Evariste Galois, seorang matematikawan Perancis (1811 -- 1832)
\end{itemize}

\end{document}
